O trabalho realizado cumpriu com o proposto de demonstrar a aplica��o de uma curva fractal na miniaturiza��o de uma antena de microfita com patch retangular. Os resultados alcan�ados est�o de acordo com o observado na literatura do assunto \cite{ref4}.

Destaca-se a relev�ncia do desempenho das antenas projetadas como satisfat�rio em contraste com a utiliza��o de um m�todo
n�o preciso de projeto empregado, no caso o modelo da linha de transmiss�o.

Demonstrou-se ao longo do presente texto, os conhecimento adquiridos ao longo da gradua��o de Engenharia El�trica, sobretudo, no que concerne a �rea de Telecomunica��es mais especificadamente � �rea de Propaga��o e Antenas que foram utilizados para compreens�o geral do assunto e a para a devida execu��o do desenvolvimento do trabalho.

Para trabalhos futuros sugere-se a aplica��o de outros contornos fractais e a constru��o das estruturas para medi��es experimentais \textit{indoor}.

Por fim, destaca-se tamb�m a relev�ncia do desenvolvimento do trabalho de conclus�o de curso como instrumento de aprimoramento
do aluno de Engenharia El�trica e como medida de sua capacidade de aplica��o do conte�do estudado durante a participa��o no projeto de inicia��o cient�fica da UFRN/CNPQ.


%%%%%%%%%%%%%%%%%%%%%%%%%%%%%%%%%%%%%%%%%%%%%%%%%%%%%
% Modelo para escrever TCCs, disserta��es e teses utilizando LaTeX, ABNTeX e BibTeX
% Autor/E-Mail: Robinson Alves Lemos/contato@robinson.mat.br/robinson.a.l@bol.com.br
% Data: 19/04/2008 
% Colaboradore(s)/E-Mail(s):
% Caso queira colaborar, entre em contato pelo e-mail e informe altera��es que realizou.
%%%%%%%%%%%%%%%%%%%%%%%%%%%%%%%%%%%%%%%%%%%%%%%%%%%%%